% Source: https://artofproblemsolving.com/community/c3h3359785_400th_post_problems
% Stable identity: aops-topic-3359785-q15
Measure speed in tens of mph, writing it as $s$. Success probability is $p_s=\min(1,\dfrac{14-s}{10})$ and reward on success is $s$. On throw one, the expected reward is $5\cdot\dfrac{9}{10}=\dfrac{9}{2}$. On throw two, speeds $6,4$ have probabilities $\dfrac{9}{10},\dfrac{1}{10}$, so the expected reward is $\dfrac{9}{10}\cdot\dfrac{4}{5}\cdot6+\dfrac{1}{10}\cdot4=\dfrac{118}{25}$. For throw three, speeds $7,5,3$ have probabilities $\dfrac{18}{25},\dfrac{7}{25},0$, respectively: the failed-speed-$4$ branch has probability zero. Thus the expected third reward is $\dfrac{18}{25}\cdot\dfrac{7}{10}\cdot7+\dfrac{7}{25}\cdot\dfrac{9}{10}\cdot5=\dfrac{1197}{250}$. The total is $\dfrac{9}{2}+\dfrac{118}{25}+\dfrac{1197}{250}=\dfrac{1751}{125}$, already reduced. Therefore $\boxed{n=125}$.
